\section{Minimal extremal windows with certified binomial bounds}
\label{sec:minimal-window-integration}

The incoming archive \path{ProveIt_Unknot_Extremal_Windows_2026-10-07.zip}
is now preserved as report 20. It adds a stronger size analysis for an alternative window schedule. The
schedule is now available in recognition through the
\code{--window-strategy} option with value \code{minimal}, and through
\code{window --minimal} for exact partial homology.
The existing allocation-pruned strategy remains the default. Both compute
exact requested ranks, but they construct different guard complexes and
have different size guarantees.

\subsection{Cancellation and truncation do not commute}
Let $K=k+1$ when the desired raw degrees are $0,\ldots,k$.
The new strategy extends the whole retained complex by one crossing,
deloops, exhaustively cancels units, and only then discards degrees above
$K$. In particular it constructs temporary degree $K+1=k+2$.
The contractible complex $X\xrightarrow{1}X$ in degrees $K,K+1$
explains the order: cancellation removes it, while early truncation leaves
a nonzero stalk in degree $K$. The latter is artificial guard homology.

This does not contradict the exactness of the earlier support-window
algorithm. That algorithm returns degrees below the guard and has a
separate bounded-suffix support proof. But that proof does not establish
the binomial size bound derived here. The two implementations are kept
explicitly separate; the new strategy never suppresses temporary objects
before cancellation. Its object ceiling therefore includes their allocation.

\subsection{Minimality controls all exhaustive pivot choices}
In the finite dotted matching category over $\F$, let $J$ consist of
all morphisms between distinct matchings and all equal-matching morphisms
with zero constant coefficient. This is the nilpotent radical; the
semisimple quotient has one scalar factor per matching. The positive
weight $m-c(a,b)+2|S|$ from the radical analysis gives a finite nilpotence
bound. An entry outside $J$ is precisely a unit between equal matchings.
Exhaustive scalar-unit cancellation therefore leaves a minimal complex,
meaning one whose differential lies in $J$.

Every bounded finite complex splits, by the invertible changes of basis
from Gaussian cancellation, into a minimal complex and contractible disks.
If minimal complexes $M,N$ are homotopy equivalent, their homotopy inverse
equations reduce modulo $J$ to degreewise inverse maps. Lift these inverses
using the finite geometric series in the nilpotent matrix ideal. The
chain maps are thus degreewise invertible and give a chain isomorphism.
Consequently the minimal matching multiplicity in every homological degree
is unique, although matrix entries and computation times can depend on
the chosen pivots.

Write $\tau_K$ for brutal upper truncation. Consider additive extension
functors $F_t$ that preserve chain homotopies and have degree amplitude
$[0,1]$. Starting from a common initial complex, compare
\[
 A_t=\tau_K\operatorname{Min}(F_t A_{t-1}),\qquad
 B_t=\tau_K R_t(F_t B_{t-1}),
\]
where $R_t$ performs any collection of valid unit cancellations. Then
\begin{equation}\label{eq:minimal-window-domination}
 B_t\simeq A_t\oplus E_t[K]
\end{equation}
for a stalk object $E_t$ in degree $K$. In particular $A_t$ has no larger
matching multiplicities than $B_t$ in any degree.

To prove this, extend the inductive decomposition. Its extra stalk extends
only into degrees $K,K+1$. Uniqueness of minimal complexes identifies the
minimal part of the extension with the sum of the desired minimal extension
and a minimal complex supported in those two degrees. The partially reduced
comparison complex is \emph{chain-isomorphic} to that sum plus contractible
disks. Truncate this chain isomorphism: disks below $K$ stay contractible,
disks above $K$ disappear, and disks crossing the cutoff become stalks in
$K$. The extra minimal two-degree piece also becomes a stalk. This proves
\eqref{eq:minimal-window-domination}. It never assumes that brutal
truncation preserves an arbitrary homotopy equivalence.

\subsection{The geometric certificate and inherited object count}
The external ingredient is
\href{https://arxiv.org/html/2601.02119v1}{Kelom\"aki--Sch\"utz,
Lemma 5.5}, checked directly against the primary paper during this review.
For their comparison scan on a nice disc order, its degree-$j$ object
count after $t$ crossings is bounded by $\binom tj$ before final closure
and by $2\binom nj$ at closure. Reducing its coefficients modulo two
preserves the allowed cancellations. Applying
\eqref{eq:minimal-window-domination} transfers these object bounds to our
exhaustive cancellation policy.

The production certificate verifies an actual crossing permutation,
absence of projection loop edges, connected processed and remaining
vertex sets, consecutive attaching slots at every intermediate crossing,
and the unique final four-slot closure. For a validated plane projection,
a cut with connected sides is a bond; its dual cycle supplies a separating
curve and a disc presentation. Consecutive attachments are checked
explicitly. The certificate records the measured maximum boundary size
$W$, not an estimate of the best possible girth.

When no order is supplied, the implementation first tries the ordinary
fast greedy order. If its certificate fails, a polynomial search attempts
a connected-prefix, connected-suffix order on the loopless two-connected
projection. It repeatedly chooses a vertex attached to the prefix whose
removal leaves the suffix connected, breaking ties by the next cut size.
The block decomposition of the connected suffix guarantees a permissible
extension under the stated two-connectedness hypothesis. Direct candidate
connectivity checks give a conservative $O(n^3)$ construction cost.
The resulting order is reverified. If construction fails, the original
order remains valid for exact computation and only the stronger complexity
certificate is withheld. Supplied orders are verified without replacement.
All planning and verification is charged to the query deadline.

At each certified stage the implementation checks the actual retained
degree counts against the binomial bound. This is a diagnostic consequence
of the proof, not a substitute for the geometric hypothesis. Tests also
compare matching-by-degree profiles across standard, residue, and adaptive
exhaustive cancellation. Optional dense coefficient operations preserve
the same exact category and can be used as well.

\subsection{A uniform bound without the all-state-circle factor}
Put
\[
 B(n,K)=\sum_{j=0}^{\min(n,K)}\binom nj.
\]
Retained population is at most $2B(n,K)$. One crossing creates at most
eight delooped children per object, so even the fully allocated temporary
extension has population $O(B(n,K))$. There are quadratically many
possible entries and linearly many pivots; a pivot changes quadratically
many entries. Each coefficient computation costs
$\operatorname{poly}(n,W)2^{O(W)}$ bit operations. Thus the new strategy
has the bounds
\[
 T\le\operatorname{poly}(n,W)2^{O(W)}B(n,k+1)^3,
 \qquad
 S\le\operatorname{poly}(n,W)2^{O(W)}B(n,k+1)^2.
\]
The implementation disables the cross-stage shape cache for this strategy
and caps its per-stage coefficient-composition memo at 4096 entries.
This avoids silently charging cubic memo storage as quadratic matrix
storage. Disc matching and plan metadata has $2^{O(W)}$ possible types
per stage; accumulated stage labels add polynomial factors.

For $k=O(\log(n+2))$ and $W=O(\log^2(n+2))$, these bounds are
quasi-polynomial. Unlike the previous elementary support-window estimate,
this bound has no all-zero smoothing circle factor. It requires a
certified nice order and this particular cancellation schedule. It is
not asserted for early-allocation pruning or arbitrary uncertified
frontiers.

The API \code{khovanov\_minimal\_window\_auto} accepts an arbitrary raw
interval $[a,b]$. On one side it computes through upper degree $b$ and
returns only the requested interval. Mirroring instead computes through
$n-a$ and maps degrees back by $h\mapsto n-h$. The default mirror choice
minimizes this raw upper depth; it does not estimate wall time.
The complexity parameter is the selected upper depth, not the width
$b-a$. The \code{binomial\_bound\_certified} field records whether the
nice-order premise was established. A failed certificate does not weaken
the exactness of the returned queried ranks.

Recognition can use this strategy inside the existing bounded widening
portfolio. Its same normalized-band comparison, local allowance, global
limits, restart rule, and complete fallback remain in force. Agreement
with the unknot on a partial band is still inconclusive. Widening restarts;
it does not attempt to reconstruct discarded chain groups. No new default
policy or general quasi-polynomial verdict guarantee is introduced.

\subsection{Raw endpoint depth has an exact padding obstruction}
Insert $s$ cancelling positive-negative crossing pairs into a diagram.
Its knot type and normalized Khovanov homology do not change, while its
crossing count rises by $2s$ and its negative PD crossing count by $s$.
Therefore the raw profile obeys
\[
 \beta_{D_s}(r)=\beta_D(r-s).
\]
Every discrepant degree relative to the unknot profile shifts by $s$,
and its distance from either raw endpoint increases by $s$. Such padding
can be inserted with constant additional frontier width. Consequently
bounded width alone does not force a logarithmic detecting endpoint depth.
Ordinary diagram simplification removes this particular padding, so this
is an obstruction to an unqualified raw-window theorem, not a lower bound
for the whole preprocessing pipeline. A theorem covering all reduced
inputs would need additional information. Even shallow agreement on an
unknot needs a positive certificate or eventual full coverage.

\subsection{Validation and comparison with the current window engine}
The integrated suite passes 277 tests. New checks compare 100 random knot
braids and arbitrary intervals with complete production homology, using
three exhaustive reducers and forced dense composition on the first 15
inputs. They check intermediate $d^2=0$, certified binomial inequalities,
mirror labels, invalid orders, padding shifts, time and allocation limits,
and the guard-disk counterexample. The CLI and recognition portfolio are
tested separately, including an agreeing partial query followed by exact
fallback. The actual stage profiles are also compared across reducers.

The archive's 14 tests pass locally. Its separately rerun independent
cube audit compares 393 windows on 100 diagrams, checks 1772 certified
stages and 15 padding cases, and reports no failures. Its archived adapter
is not applied: the production implementation here targets the current
array-based \code{FastScan}, with its existing coefficient engines and
adaptive cancellation. Source archive files remain unchanged; review
runs use a scratch extraction. Local logs are retained under
\path{synthesis/data/}.

\begin{center}
\begin{tabular}{@{}lrrrr@{}}
\toprule
Input & Full/support & Full/minimal & Support/minimal & A/A \\
\midrule
weaving\_14 & 4.52 & 4.09 & 0.895 & 1.002 \\
torus\_3\_4 & 1.67 & 1.29 & 0.777 & 0.996 \\
padded\_trefoil & 1.74 & 1.20 & 0.708 & 1.009 \\
conway & 1.63 & 1.35 & 0.834 & 0.994 \\
kinoshita\_terasaka & 1.72 & 1.43 & 0.831 & 1.017 \\
stress\_braid5\_36 & 425.47 & 323.49 & 0.792 & 0.999 \\
\bottomrule
\end{tabular}
\end{center}
Ratios are medians of paired times; larger than one favors the denominator.


These seven-round paired measurements use common supplied crossing orders
and depth-two raw queries. Full homology and a window have different
output contracts; both window strategies answer exactly the same interval.
A/A compares identical full-homology calls. Order construction and imports
are excluded, while the minimal strategy's certificate verification is
included. These are backend query measurements, not default recognition
speedups. In particular, a supplied nice order can differ from the ordinary
recognizer's preferred order.

The new strategy is slower than allocation pruning on all six cases.
It sometimes saves retained guard objects nevertheless: on the 14-crossing
weaving case, peak pre-cancellation allocation is 492 rather than 516,
versus 5046 for full homology in that order. On Conway it instead allocates
84 objects versus 42 for the support strategy, because it constructs the
temporary degree needed for cancellation. Its practical tradeoff is an
additional proved size guarantee under explicit hypotheses, not a measured
replacement for the faster default on this corpus.
